% Newton's method for unconstrained minimization: at each iterate fit a
% quadratic from the gradient and Hessian, head toward its stationary point
% (or downhill if that direction is not a descent direction), and choose the
% step with a line search started from the natural Newton step alpha = 1.
% Symbols follow Alg. 4.6 in Martins & Ning, Engineering Design Optimization
% (mdobook.github.io), with Chapra's E_s for the tolerance.
% Cropped to its own bounding box so it can be dropped onto a Xournal++ page
% as an image.
%
%   latexmk -pdf newton-pseudocode.tex
\documentclass[border=6pt,varwidth=12cm]{standalone}

\usepackage{algorithm2e}
\usepackage{amsmath}

\RestyleAlgo{plain}          % no float, no ruled box
\DontPrintSemicolon
\SetKwComment{Comment}{// }{}
\SetKwInOut{Input}{input}
\SetKwInOut{Output}{output}

\begin{document}
\begin{minipage}{\linewidth}
\begin{algorithm}[H]
  \textbf{Newton's Method}
  \BlankLine
  \Input{starting point $\vec{x}_0$, tolerance $E_s$}
  \Output{$\vec{x}^*$, estimate of a local minimizer}
  \BlankLine
  $k \leftarrow 0$\;
  \While{$\|\nabla f(\vec{x}_k)\| > E_s$}{
    $\vec{p}_k \leftarrow -H(\vec{x}_k) \backslash \nabla f(\vec{x}_k)$ \tcp*{stationary point of quadratic fit}
    \If{$\nabla f(\vec{x}_k)^\mathsf{T} \vec{p}_k \ge 0$}{
      $\vec{p}_k \leftarrow -\nabla f(\vec{x}_k)$ \tcp*{not downhill: use steepest descent}
    }
    $\alpha_k \leftarrow$ line search along $\vec{p}_k$ with $\alpha_\text{init} = 1$\;
    $\vec{x}_{k+1} \leftarrow \vec{x}_k + \alpha_k\, \vec{p}_k$\;
    $k \leftarrow k + 1$\;
  }
  \Return $\vec{x}^* = \vec{x}_k$\;
\end{algorithm}
\end{minipage}
\end{document}
